\section{Problemas de Diseño de Redes}

\subsection{Introducción}

Los problemas de diseño de redes constituyen un área de problemas donde se tiene que encontrar la configuración o el diseño
óptimo de una red, de forma que se cumplan ciertos criterios de servicio y conectividad, dependiendo del problema concreto.
Este área tiene muchas aplicaciones prácticas, fundamentalmente en problemas de redes de comunicaciones y en redes logísticas
de distribución. Cuando el objetivo no es el diseño de la red sino la optimización del uso de la misma se tienen los problemas
de optimización de flujo en red. En la siguiente tabla se muestran ejemplos de problemas de cada tipo

\begin{table}[h!]
    \centering
    \begin{tabular}{l|l}
        \textbf{Problemas de optimización del uso} & \textbf{Problemas de diseño} \\ \hline
        Problemas de transporte & Problemas de localización con costos fijos \\
        Flujo en redes de un producto & Transporte con costos fijos \\
        Flujo en redes multiproducto & Flujo en redes con costes fijos \\
        Transporte en dos etapas & Sistema de distribución en dos etapas
    \end{tabular}
\end{table}

\subsection{Sistemas de distribución}

Los problemas de planificación de la distribución de productos se clasifican en tres tipos, problemas estratégicos (plazo
largo, incertidumbre grande), problemas tácticos (plazo medio, incertidumbre menor) y problemas operacionales (plazo muy
corto, incertidumbre pequeña). Para el diseño de un sistema de distribución de productos hay 3 esquemas o sistemas básicos

\begin{itemize}
    \item \textbf{Transporte directo}. Los productos se transportan directamente desde las plantas hasta los puntos de
    distribución finales. El problema de diseño para este tipo de sistema para un único producto se modela como un problema
    de localización con costes fijos en una de sus tres variantes, UFLP (Uncapacited Facility Location Problem), CFLP
    (Capacited Facility Location Problem) o SSCFLP (Single Source Capacitated Facility Location Problem). La optimización de
    un sistema ya diseñado se reduce al problema de transporte clásico.
    \item \textbf{Almacenamiento}. Es el método tradicional de almacenar los productos en tanques/almacenes intermedios. El problema de
    optimización para un producto es un problema de transporte en dos etapas que tambien puede modelizarse como problema de
    flujo en redes.
    \item \textbf{Terminales de carga (Crossdocking)}. Este método se conoce tamién como distribución just-in-time. Los
    productos se envían a una instalación ``crossdock'' con terminales de carga donde son ordenados, las cargas se consolidan
    con otros productos y se envian con trailers a los puntos de demanda sin necesidad de almacenarlos. El problema de
    optimización de un sistema con terminales de carga es un problema de flujos en redes multiproducto, y el problema de diseño
    es un problema de flujos en redes multiproductos con costos fijos.
\end{itemize}

\subsection{Problema de Flujo en Redes con Costos Fijos (FCNP)}

Se tiene un grafo dirigido $G=(N,A)$ con $N=\{1,\dots,n\}$ conjunto de nodos y $A$ un conjunto de arcos, donde cada arco
corresponde a un par $(i,j)$ tal que $i,j\in N$. Asociado a cada arco hay una variable de flujo con valores 
$0\leq x_{ij}\leq u_{ij}$ donde $u_{ij}$ es una cota superior del flujo en el arco $(i,j)$ y una variable binaria 
$y_{ij}\in\{0,1\}$ que representa el ``uso'' del arco en la solución final (cuando $x_{ij}>0$).Para ello es necesaria la
restricción
\[
    x_{ij}\leq u_{ij}y_{ij}
\]
Las variables $y_{ij}$ son las variables de diseño y tienen asociadas en la función objetivo un costo fijo $f_{ij}$. En 
ocasiones las variables $x_{ij}$ tienen asociado un costo variable $c_{ij}$. Finalmente, $b_i$ es la oferta o demanda del
nodo $i$. Se formula el modelo, suponiendo la red balanceada, como

\begin{alignat}{3}
    \text{minimizar:}   & \sum_{(i,j)\in A}(f_{ij}y_{ij} + c_{ij}x_{ij})    &                       \\
    \text{sujeto a:}    & \sum_{j=1}^nx_{ij}-\sum_{j}x_{ji}=b_i                 & \forall i\in N        \\
                        & x_{ij}\leq u_{ij}y_{ij}                           & \forall (i,j)\in A    \\
                        & x_{ij}\geq 0                                      & \forall (i,j)\in A    \\
                        & y_{ij}\in\{0,1\}                                  & \forall (i,j)\in A
\end{alignat}

\subsection{Problema Conector}

El Problema Conector consiste en encontar la forma de conectar $n$ puntos (nodos) entre si con costo o distancia total mínima 
a partir de un árbol $G=(N,A)$ donde cada arista $(i,j)$ tiene asociado un costo $\omega_{ij}$. Un subgrafo expandido conexo 
es un grafo $G'=(N,T)$, con conjunto de aristas $T\subseteq A$ que conecta entre si todos los nodos de $N$, y que tiene 
asociado un costo total $\omega(T)=\sum_{(i,j)\in T}\omega_{ij}$. El problema conector se reduce a encontrar el subgrafo 
expandido conexo de menor costo total.

\subsubsection{Minimum Spaning Tree (MST)}

Un árbol es un grafo conexo y sin ciclos. En la mayoria de problemas prácticos de conexión se puede suponer que 
$\omega_{ij}>0$ para toda arista $(i,j)\in A$. En este caso el Problema Conector tiene como solución un árbol expandido,
pues si existiese un grafo expandido con algún ciclo simplemente se puede eliminar una arista manteniendo la conexión entre
todos los nodos y reduciendo \textbf{siempre} el costo total. 

Por tanto el Problema Conector se ve reducido al problema MST (Minimun Spanning Tree), que puede ser formulado como Problema 
de Flujo en Redes con Costos Fijos. Primero se duplica cada arísta en dos arcos, pues el Problema de Flujo en Redes con Costos
Fijos es dirigido, $i\to j$ y $j\to i$. Si se van a conectar $n$ puntos, se elige uno arbitariamente (nodo $k$) y se toman 
ofertas/demandas de $b_k=n-1$ y $b_i=-1$ para $k\neq i$. Estas ofertas y demandas obligan a conectar todos los puntos 
proporcionando una solución de grafo conexo. En la función objetivo solo aparecerán los costos de los arcos por las variables 
de diseño binarias $y_{ij}$ (no hay costos variables). Por último, para asociar las variables 
$x_{ij}\geq 0$ y $y_{ij}\in\{0,1\}$ en las restricciones $x_{ij}\leq u_{ij}y_{ij}$ se puede usar como $u_{ij}=n-1$ para todo 
arco $(i,j)\in A$.

\vspace*{2mm}

A pesar de que el problema MST puede ser formulado como problema mixto de redes con costes fijos, hay algoritmos que lo
resuelven fácilmente sin necesidad de plantear el modelo.

\newpage

\subsubsection{Algoritmo de Kruskal}

Es un caso particular de los algoritmos ``greedy'' pues en cada paso se toma la arista o variable de costo mínimo, y no se
vuelve a reoptimizar o corregir. Para el MST el algoritmo es el siguiente

\begin{itemize}
    \item \textbf{Etapa inicial}. Se elige una arista de costo mínimo
    \item \textbf{Etapa general}. Si se han elegido ya aristas $A'=\{a_1,\dots,a_i\}$ se elige la siguiente arista 
    $a_{1+1}\in A-A'$ de forma que se verifique que en $A'=\{a_1,\dots,a_i,a_{i+1}\}$ no hay ciclos y el costo 
    $\omega(a_{i+1})$ es mínimo.
    \item \textbf{Etapa final}. Se termina cuando ya no se pueda realizar la etapa anterior (siempre en $n-1$ etapas).
\end{itemize}

\subsubsection{Algoritmo de Prim}

Se trata de otra implementación del algoritmo greedy. En vez de ir añadiendo aristas hasta formal el árbol final (con las 
comprobaciones de ciclos que conlleva), en este método se va aumentando un árbol expandido $T_j$ desde uno con un único nodo y
ninguna arista hasta uno con $n$ nodos y $n-1$ aristas.

\begin{itemize}
    \item \textbf{Etapa inicial}. $T_0$ consiste de un único nodo.
    \item \textbf{Etapa general}. Dado $T_{j-1}$ árbol conteniendo $j$ nodos y $j-1$ aristas, se añadie la arista de costo o
    peso mínimo entre todas las aristas con un extremo en $T_{j-1}$ y otro en $T_{j-1}'$.
    \item \textbf{Etapa final}. Se termina con el árbol $T_{n-1}$ (siempre en $n-1$ etapas).
\end{itemize}

\subsubsection{Capacitated Minimum Spaning Tree (CMST)}

El Problema del Mínimo Árbol con Capacidades (CMST) es una versión del problema MST con la concición adicional de que ninguno
de los sub-árboles saliendo desde la raíz tenga más de $k$ nodos. Salvo para los casos triviales de $k=1$ (árbol estrella) y 
$k\geq n-1$ (MST) se trata de un problema mucho más dificil de resolver que el MST y no existe ningún algoritmo simple
(polinómico) pàra su resolución.


En otra versión del problema existen demandas $d_i$ desde los nodos $i\neq 1$ al nodo raíz o desde el nodo raíz al resto de
nodos. En este caso la restricción de capacidad es un límite $C$ sobre la capacidad de tráfico que puede manejar. Entonces se
debe cumplir la restricción $\sum_{i\in T_k}d_i \leq C$ donde $T_k$ es el conjunto de nodos del árbol conectado a cualquier
puerto $k$.

\vspace*{2mm}

El problema CMST se puede formular como un Problema de Flujo en Redes con Costos Fijos con las siguientes características.
\begin{enumerate}
    \item Costos fijos. $f_{ij}=\omega_{ij}$ mientras que los costos variables son nulos para todo arco $(i,j)\in A$.
    \item Ofertas/demandas. Hay dos formas de formularlas, dependiendo de si el flujo es de los nodos a la raíz $r$ o de la 
    raíz a los nodos. Si el tráfico es de los nodos a la raíz, en el modelo de redes se toman como ofertas/demandas $b_i=1$
    para $i\neq r$ y $b_r=-(n-1)$, y para el caso de demandas de tráfico $b_i=d_i$ para $i\neq r$ y $b_r=-\sum d_r$. Cuando el
    tráfico es de la raíz a los nodos se toman como ofertas/demandas $b_i=-1$ para $i\neq r$ y $b_r=(n-1)$, y para el caso de
    demandas $b_i=-d_i$ para $i\neq r$ y $b_r=\sum d_r$. En todos los casos la red está balanceada.
    \item Capacidades. En el caso del tráfico de los nodos a la raíz, $u_ir=k$ y para el resto en principio no hay restricción
    de capacidad. Si el tráfico va del nodo raíz hasta el resto de nodos, $u_ri=k$ y para el resto no hay restricción de
    capacidad.
\end{enumerate}

\newpage

\subsubsection{Métodos heurísticos CMST}

Existe uma variedad de métodos heurísticos para el problema CMSP. Heurísticas de construcción conocidas son la heurística de
Eseau-Williams, que tiene su origen en el conocido algoritmo de Clarke-Wright para el problema de rutas con capacidades y se
basa en el concepto de ahorro (saving). Otra más sencilla es el algoritmo de Prim modificado, el método consiste en ir
construyendo subárboles con raíz en el nodo base $r$, con el criterio del método greedy de Prim y respetando la restricción
de capacidad (se puede realizar de forma secuencial o en paralelo).

\subsubsection{Problema de Steiner}

Otro Problema Conector que surge en la práctica es el Problema de Steiner. En este problema se tiene, al igual que en el MST,
un grafo $G=(N,A)$ con pesos o distancias en las aristas $\omega_{ij}$ para $(i,j)\in A$ y uin subconjunto de nodos
$T\subseteq N$ conocidos como terminales. El problema consiste en encontrar un subgrafo que conecte los nodos terminales entre
sí y que tenga peso total mínimo. Si los pesos son positivos la solución final tendrá forma de árbol, conocido como árbol de
Steiner. Si $T=N$ el problema se reduce al MST, mientras que si $|T|=2$ se reduce al del camino más corto entre los dos nodos
en $T$. Para el resto de casos, el problema de Steiner es un problema dificil para el que no existe ningún algoritmo
polinómico.

El problema del mínimo árbol de Steiner se puede formular como un problema de flujo en redes con costos fijos con las siguientes
características

\begin{enumerate}
    \item Se transforma el grafo en uno dirigido duplicando cada arista $(i,j)$ en dos arcos dirigidos $(i,j)$ y $(j,i)$.
    \item Costos fijos. $f_{ij}=\omega_{ij}$ mientras que los costos variables son nulos para todo arco $(i,j)\in A$.
    \item Ofertas/demandas. Se escoge arbitrariamente un nodo terminal como raíz $r\in T$ y se toman como ofertas/demandas
    $b_r=|T|-1$ y $b_i=-1$ para $i\in T$, $i\neq r$ y $b_i=0$ para todo $i\notin T$.
    \item Capacidades. Las capacidades de los arcos se toman igual al flujo total, que es $|T|-1$.
\end{enumerate}

\subsection{Modelos de Diseño Redes Multiproducto con Costos Fijos}

Los problemas de diseño de redes multiproducto son una extensión a varios productos de los problemas de diseño vistos
anteriormente. Los elementos básicos son un grafo dirigido $G=(N,A)$ donde $N$ es el conjunto de nodos y $A$ el conjunto de
arcos, y un conjunto de productos $K$. En los siguientes modelos tenemos de nuevo dos clases de variables de decisión, la
cantidad de flujo del producto $k$ que pasa por el arco $(i,j)$ que denotaremos como $x_{ij}^k$ y $y_{ij}$ que indican si el
arco $(i,j)$ se usa o no en el diseño solución. Para cada $x_{ij}^k$ tendremos un coste unitario $c_{ij}^k$ asociado y en el
caso de las $y_{ij}$ un coste fijo $f_{ij}$. También hay una serie de ofertas/demandas $b_i^k$ para cada nodo $i\in N$.

\begin{flushright}
    Hasta aquí llegué, consultar más en los apuntes que me va a dar un chungazo de escribír\footnote{Es mucho y repetitivo de 
    anteriores temas pero generalizado a $K$ productos, consultar en el document ``Diseño redes multiproducto.pdf. Más adelante
    esta sección será completada correctamente. La formula objetivo y las restricciones son bastante sencillas de intuir.''}
\end{flushright}